\section{Decision-scale qualification: proofs}\label{sec:proofs96}
\begin{proof}[Proof of Proposition~\ref{prop:quadrature96}]
For every real threshold $\eta$, the function $\phi_\eta=(v-\eta)_+$ is convex and $K$-Lipschitz. Conditional Jensen and the chord inequality give
$M_n\phi_\eta\leq P\phi_\eta\leq E_n\phi_\eta$.
Apply the AVaR threshold formula
$\inf_\eta\{\eta+(1-\alpha)^{-1}\mu\phi_\eta\}$
for each of the three probability measures. Taking infima preserves the two orders and proves the enclosure. In particular, no interchange of an action infimum and a risk supremum occurs.

Here is a uniform bound which remains valid at kinks. Let $h=2R/n$. Choose one-sided slopes of $\phi_\eta$ at each bin's outer endpoints. The slope increases across each bin; its total increase over the whole interval is at most $2K$. The difference between the endpoint average and midpoint value on bin $i$ is no greater than $h/4$ times the slope increase across that bin. This follows by integrating the left and right slopes over the two half-bins. Choose the left endpoint right derivative and right endpoint left derivative; jumps at shared endpoints may be omitted, which only decreases the sum of within-bin slope increases. Therefore
\[
 0\leq E_n\phi_\eta-M_n\phi_\eta
       \leq\frac{h}{4n}(2K)=\frac{KR}{n^2}.
\]
For completeness, the half-bin expression equals one half the integral of the difference between the right and left slopes at equal distances from the midpoint, so its upper bound also follows directly without assuming a classical second derivative. The preceding uniform bound passes through the two infima and gives \eqref{eq:quad96}. Monotonicity and cash invariance add $e_-+e_+$ for endpoint error. They also charge the actual terminal controller's one-sided error only on the upper side.

For a product law, apply the one-coordinate midpoint--endpoint comparison successively, holding the other coordinates fixed. Products of positive averaging operators preserve inequalities and telescope. Each partially averaged positive-part function remains convex and has the same Lipschitz modulus in each unintegrated coordinate. Summing the one-coordinate error bounds and then taking the threshold infimum proves the product statement. This argument uses neither a discrete economic transition assumption nor differentiability of $v$.
\end{proof}

\begin{proof}[Proof of Theorem~\ref{thm:power96}]
Every service's regret belongs to $[0,\tau]$ at every possible initial state. The difference of two complete-policy values therefore belongs to $[-\tau,\tau]$, because the common original optimal value cancels. Each interval intersection retains the true paired difference. A difference of two intervals of width at most $w$ has width at most $2w$. Intersecting with a valid support cannot increase that width.

Hoeffding's inequality for a variable of range length $2\tau$ gives
\[
 \Pr\{|\overline D-\mathbb ED|>r\}
       \leq2\exp\{-Nr^2/(2\tau^2)\}.
\]
A union bound across the $H$ fixed contrasts gives the stated coverage. Pairwise contrasts need not be mutually independent; only the initial states within each contrast must be independent, conditional on the frozen preparation. Replacing the unobserved sample mean by its numerical lower and upper bounds yields \eqref{eq:paired96}. Their difference is at most $2(r+w)$. This is an application of the established concentration inequality \citep{hoeffding1963}, not a claim to a new concentration principle.

For the chosen rational parameters, $Nr^2/(2\tau^2)=8$. The inequality $24e^{-8}<1/20$ can be checked entirely with rational arithmetic: the partial sum $\sum_{j=0}^{11}8^j/j!$ is greater than $480$. Equation~\eqref{eq:power96} is an exact rational comparison. In numerical classification the margin is rounded inward and the reported confidence endpoints outward. No observed range or variance is used to choose the radius.
\end{proof}

\paragraph{Totalization of the admitted service.}
The first-date objective with the original $V_1$ has action Lipschitz constant less than $7$: the primitive action derivative is at most $3/4$, and the discounted continuation contribution is at most $(15/16)(15/d)(3d/8)$. The feasible interval has length at most $1/4$. A uniform net with $m=1024$ intervals therefore gives an all-action lower value by subtracting $7/(8m)$ from the smallest grid cost lower bound. The grid point with the smallest upper value supplies a feasible policy upper value.

The exact fallback uses $n=256$ innovation bins and the preceding convex-order rule. Exact one-date rational supports have gaps at most $2^{-30}$; the terminal controller contributes $2^{-20}$. The resulting quadrature width, including these allowances, is below $1/8192$ for both risk levels. Adding $7/(8m)$ leaves the original-optimum bound below $1/256$. Downward rounding of the first action to 40 bits costs at most $7\,2^{-40}$ on each side and fits strictly inside both remaining allowances. All values, risk weights, feasible actions and endpoint computations in this final construction are rational, and the terminal strong-convexity bisection reaches its stated gap in finitely many steps. Thus the fallback defines a finite valid action at every dyadic initial observation in the target distribution. For arbitrary real initial states, the separate acquisition transfer remains required; the confirmation theorem here concerns the declared dyadic initial law.

The floating implementation may reach the two targets much sooner. That does not permit removing the exact fallback from the policy definition, treating rejected proposals as missing observations, or asserting that the fallback was exercised when the actual counter is zero. Hardware or resource failure aborts the experiment and is recorded as a failure, not a change in its conditional sampling law.

\begin{proof}[Proof of Proposition~\ref{prop:purchase96}]
Subtract the two complete acquisition-and-deployment accounts. Using the upper policy, acquisition and fee bounds and the lower computational saving bounds above \eqref{eq:purchase96} by
$m(d^+-\lambda s^-)+\lambda k^++F^+$.
Nonpositivity gives \eqref{eq:adopt96}. Solving this scalar inequality proves the integer threshold and its sign qualifications. The separate quality inequality depends on policy cost rather than price and cannot be deduced by increasing $\lambda$. A common audit cost cancels only when it is actually common and purchased in both alternatives; otherwise its incremental charge remains in $k^+$ or $F^+$.
\end{proof}

\begin{proof}[Proof of Corollary~\ref{cor:sharpness96}]
Write $h=2R/n$ and $k=n(1-\alpha)$. The two roots of $v_n-\eta_n$, where $\eta_n=K\alpha R$, are $c_n-\alpha R$ and $c_n+\alpha R$. They lie strictly inside the domain and at two distinct bin midpoints, since $\alpha n$ is an even integer. Under $M_n$ the mass strictly above $\eta_n$ is $(k-1)/n$ and the mass at $\eta_n$ is $2/n$. Under $E_n$ there is no mass at $\eta_n$, and the mass above it is exactly $k/n$. Hence $\eta_n$ minimizes the AVaR threshold objective under both measures, by its left and right slope conditions.

The hinge $\phi=(v_n-\eta_n)_+$ is affine on every bin except the two whose midpoints are its roots. On either exceptional bin, the endpoint average minus the midpoint value is $Kh/4$. Each bin has weight $1/n$. The two threshold minima therefore differ by
\[
 (1-\alpha)^{-1}(E_n\phi-M_n\phi)
 =\frac{1}{1-\alpha}\frac{2}{n}\frac{Kh}{4}
 =\frac{KR}{(1-\alpha)n^2}.
\]
All equalities use the endpoint mixture's half weights at the outer boundaries. This explicit sequence establishes sharpness of the stated enclosure on its function class; it does not assert that the investment value itself equals $v_n$.
\end{proof}

